\section{What this closes and does not close}

\begin{description}[style=nextline]
\item[\texttt{SEAM.STATE.RPMIXING.01}, continuum leg.]
For an abstract holomorphic conformal net satisfying the hypotheses,
the rotational KMS extension theorem is complete: there is one
extension, its expectation is canonical, and both local
relative-commutant mixing and thermal-sector mixing are excluded.
The finite \texttt{v995} computation is exactly consistent with this
result.  Applying it to the TFPT collar remains conditional on the
continuum-net identification in the box below.  The ledger's
continuum leg therefore stays open.

\item[\texttt{ALPHA.QUILLEN.EXACT.01}.]
The theorem supplies the rigidity premise: once the U(1)
determinant-line functional is identified with the invariant modular
response of this inclusion, there is no second KMS state or
nontrivial relative-commutant direction that can alter its sector
weights.  It does not prove the Bismut--Freed/Quillen identification,
the zeta-regularized variation, or the compensation normalization.
That contract remains open.

\item[\texttt{AX.P1.01}.]
The arithmetic identity
$1/(4\cdot2\pi)=1/(8\pi)$ and the equidistribution implication are
proved under explicit operator and geometric hypotheses.
The identification of that response with the primitive seam
coefficient is not.  Thus P1 stays typed as an axiom and gains a
theorem-conditional reading, not a marker upgrade.
\end{description}

\section{Remaining gap}

\begin{center}
\setlength{\fboxsep}{10pt}
\fcolorbox{red!65!black}{red!3}{
\begin{minipage}{0.91\textwidth}
\textbf{Remaining gap (exact statement).}
For the TFPT $M=1$ QWZ collar with its four quarter-twisted deck
labels, prove that the scaling-limit local von Neumann algebras are
the hyperfinite type-$\mathrm{III}_1$ factors of a conformal net and
identify, as conformal nets with vacuum, common stress tensor,
rotation dynamics, and $\Z_4$ action, the resulting concrete
inclusion with
\[
 \A=\B^{\Z_4}\subset\B,\qquad
 [\B:\A]=4,\qquad
 \A(I)'\cap\B(I)=\C,\qquad
 \mu_{\B}=1,\quad c_{\B}=8.
\]
The identification must construct the non-bilinear spin/twist
fields, verify relative locality and the commutative index-four
Q-system, and show that the extension is the holomorphic $(E_8)_1$
net rather than merely matching its finite covariance, fusion, and
index arithmetic.  For the P1/Quillen application, the same
identification must additionally show that the normalized
determinant-line modular response is the additive transitive
four-channel response of Lemma~\ref{lem:equidistribution}, with
$\beta_{\mathrm{angle}}=2\pi$.

\medskip
This is the MMST/Osborne--Stottmeister leg beyond fermion bilinears:
the existing theorem controls the bilinear
$\mathfrak{so}(16)_1$ currents, but not the spinor extension fields
that generate $(E_8)_1$.  Finite factoriality
$Z(M_{256})=\C$, the exact covariance \eqref{eq:frozen-C}, the
$33\to0$ KMS rank test, and the $196608\to0$ invariant-extension
test do not prove this type-$\mathrm{III}_1$ scaling-limit
identification.
\end{minipage}}
\end{center}

\section*{Conclusion}
\addcontentsline{toc}{section}{Conclusion}

The abstract operator-algebraic theorem is closed at its natural
strength.  Existence comes from the vacuum-sector rotational Gibbs
state of the holomorphic extension.  The index-four group average is
the unique expectation by the finite-index Radon--Nikodym theorem
and irreducibility.  Takesaki's theorem gives the explicit modular
restriction and intertwining chain.  Gauge-invariant extension
uniqueness is an elementary averaging identity, while unconditional
rotational KMS uniqueness requires both KLM ($\mu=1$ leaves only the
vacuum sector) and Longo--Tanimoto (rotational KMS extreme points are
sector Gibbs states).  Araki's theorem identifies the exact frozen
finite covariance and explains the finite KMS rank collapse.

The remaining issue is not another uniqueness argument.  It is the
construction and identification of the concrete TFPT collar
scaling-limit net, including its non-bilinear extension fields and
its determinant-line response.  That issue is isolated in the
single box above and lies outside Theorem~\ref{thm:main}.

\begin{thebibliography}{99}

\bibitem{Araki1971}
H.~Araki,
\emph{On quasifree states of CAR and Bogoliubov automorphisms},
Publ.\ Res.\ Inst.\ Math.\ Sci.\ \textbf{6} (1970/71), 385--442,
\href{https://doi.org/10.2977/prims/1195193913}
{doi:10.2977/prims/1195193913}.

\bibitem{Boeckenhauer1996}
J.~B\"ockenhauer,
\emph{An algebraic formulation of level one Wess--Zumino--Witten
models},
Rev.\ Math.\ Phys.\ \textbf{8} (1996), 925--948.

\bibitem{BoeckenhauerEvans1998}
J.~B\"ockenhauer and D.~E.~Evans,
\emph{Modular invariants, graphs and $\alpha$-induction for nets of
subfactors I},
Commun.\ Math.\ Phys.\ \textbf{197} (1998), 361--386.

\bibitem{Hiai1991}
F.~Hiai,
\emph{Minimum index for subfactors and entropy II},
J.\ Math.\ Soc.\ Japan \textbf{43} (1991), 347--379.

\bibitem{KLM2001}
Y.~Kawahigashi, R.~Longo and M.~M\"uger,
\emph{Multi-interval subfactors and modularity of representations
in conformal field theory},
Commun.\ Math.\ Phys.\ \textbf{219} (2001), 631--669,
\href{https://doi.org/10.1007/PL00005565}
{doi:10.1007/PL00005565}.

\bibitem{Longo1989}
R.~Longo,
\emph{Index of subfactors and statistics of quantum fields I},
Commun.\ Math.\ Phys.\ \textbf{126} (1989), 217--247.

\bibitem{LongoRehren1995}
R.~Longo and K.-H.~Rehren,
\emph{Nets of subfactors},
Rev.\ Math.\ Phys.\ \textbf{7} (1995), 567--597.

\bibitem{LongoTanimoto2018}
R.~Longo and Y.~Tanimoto,
\emph{Rotational KMS states and type I conformal nets},
Commun.\ Math.\ Phys.\ \textbf{357} (2018), 249--266,
\href{https://doi.org/10.1007/s00220-017-2969-8}
{doi:10.1007/s00220-017-2969-8}.

\bibitem{OsborneStottmeister2023}
T.~J.~Osborne and A.~Stottmeister,
\emph{Conformal field theory from lattice fermions},
Commun.\ Math.\ Phys.\ \textbf{398} (2023), 219--283,
\href{https://doi.org/10.1007/s00220-022-04519-4}
{doi:10.1007/s00220-022-04519-4}.

\bibitem{Takesaki1972}
M.~Takesaki,
\emph{Conditional expectations in von Neumann algebras},
J.\ Funct.\ Anal.\ \textbf{9} (1972), 306--321.

\bibitem{Takesaki2003}
M.~Takesaki,
\emph{Theory of Operator Algebras II},
Encyclopaedia of Mathematical Sciences \textbf{125},
Springer, Berlin, 2003.

\end{thebibliography}

\end{document}
